\chapter{Conclusion}
\label{ch:conclusion}

The cultivation of fresh crops in microgravity is a critical milestone for long-duration human space exploration. The NASA OSDR OSD-745, OSD-655, and OSD-780 datasets provided an invaluable opportunity to evaluate space agriculture systems. 

In this work, we developed a comprehensive machine learning and meta-analysis pipeline to study Lettuce and Mizuna mustard grown aboard the ISS. Our multivariate approach, driven by Random Forest classification and interpreted via SHAP feature importance, successfully demonstrated that space-grown leafy greens exhibit a conserved multielemental fingerprint. Crucially, while these physiological shifts were significant, they did not compromise the safety or the overall nutritional value of the crops, reaffirming the viability of the Veggie platform.

Furthermore, we established that the tabular foundation model TabPFN (Nature 2025) is exceptionally well-suited for small space biology datasets, achieving 89.6\% cross-validated accuracy and showing strong cross-species generalizability (83.3\%) under LOCOCV. This validates the existence of a universal spaceflight physiological signature across plant families.

Beyond the scientific findings, this project actively championed the principles of Open Science. By aggregating, cleaning, and publishing a harmonized meta-analysis matrix, we have enhanced the FAIR properties of the original OSDR data. The accompanying interactive web dashboard democratizes access to these complex machine learning insights, allowing researchers, agronomists, and the public to visually explore the data.

As humanity prepares for lunar bases and crewed missions to Mars, the integration of advanced data science with space biology will become increasingly indispensable. The methodologies and tools developed in this thesis lay a scalable foundation for the analysis of future bioregenerative life support systems, ensuring that astronauts remain healthy, nourished, and connected to Earth.
